Fix \(\mathbf M\in\mathcal C(\mathcal P)\) and a checked derivation \(\delta:A\Downarrow F\) in domain \(s\).  We first prove
\[
\llbracket F\rrbracket_{\mathcal P}=P_{M^s}(A)
\tag{10}
\]
by structural induction on \(\delta\).  At a source-cell leaf, \(\mathrm{def:compatibility\text{-}fiber}\) and its typed-menu realizability requirement give
\[
\llbracket\operatorname{CellVal}(s,a,v)\rrbracket_{\mathcal P}
=P_a^s(v)
=P_{M^s}(C_{a,v}).
\tag{11}
\]
The contradiction and sure-event leaves give zero and one.  If the last constructor is a disjoint sum, its exhaustive truth-table check and the induction hypotheses give
\[
P_{M^s}\!\left(\biguplus_{i\in I}A_i\right)
=\sum_{i\in I}P_{M^s}(A_i)
=\sum_{i\in I}\llbracket F_i\rrbracket_{\mathcal P}.
\tag{12}
\]
Marginalization is (12) for the finite partition over every value of the marginalized variable.  A checked complement has \(A\subseteq H\), and hence
\[
P_{M^s}(H\setminus A)
=P_{M^s}(H)-P_{M^s}(A)
=\llbracket F_H-F_A\rrbracket_{\mathcal P}.
\tag{13}
\]
For a product node, the registry checker verifies measurability with respect to pairwise disjoint sets of original district blocks; \(\mathrm{ass:district\text{-}factorization}\) and the induction hypotheses then give the product equality (6).  For a conditional node with numerator \(N\subseteq H\), its separate context certificate is accepted before the quotient is formed.  Once the certificate soundness proved below gives \(P_{M^s}(H)>0\), the induction hypotheses give
\[
\frac{\llbracket F_N\rrbracket_{\mathcal P}}
     {\llbracket F_H\rrbracket_{\mathcal P}}
=\frac{P_{M^s}(N)}{P_{M^s}(H)}
=P_{M^s}(N\mid H).
\tag{14}
\]
Registry, substitution, and transitivity nodes are sound by \(\mathrm{lem:finite\text{-}fixing\text{-}calculus\text{-}soundness}\), applied to their checked finite traces.  These are all constructors of \(\mathrm{def:finite\text{-}extraction\text{-}calculus}\), so (10) follows.

We now justify the positivity step used in (14), without treating a claimed inequality as a certificate.  The certificate grammar in \(\mathrm{def:finite\text{-}extraction\text{-}calculus}\) is well founded: a certificate for \(H\) may use only previously accepted positive denominators.  After clearing those denominators, every accepted certificate contains a rational \(\varepsilon>0\), the displayed polynomial equalities \(f_\ell(z)=0\) and inequalities \(g_j(z)\ge0\) defining its declared finite compatibility fiber, and an exact finite sign derivation of
\[
\Delta(z)\bigl(\llbracket H\rrbracket(z)-\varepsilon\bigr)\ge0,
\qquad \Delta(z)>0.
\tag{15}
\]
The primitive sign rules are equality substitution, addition and multiplication of already certified nonnegative quantities, multiplication by a rational nonnegative scalar, and the fact that an explicit square is nonnegative.  Each polynomial equality used by a rule is accepted only after fixed-order coefficient normalization.  Induction on this finite sign derivation is sound over every \(z\) satisfying the displayed constraints: equality terms vanish, each \(g_j(z)\) is nonnegative, sums and products of nonnegative real numbers are nonnegative, and squares are nonnegative.  The linear-face or Farkas form is the special case
\[
\Delta(z)\bigl(\llbracket H\rrbracket(z)-\varepsilon\bigr)
=\lambda^{\mathsf T}(Bz-d)+\gamma^{\mathsf T}z,
\qquad \gamma\ge0,quad Bz=d,quad z\ge0,
\tag{16}
\]
whose soundness follows immediately by substitution.  Direct rational certificates are the constant special case of (15).  Exact real-algebraic certificates are finite compositions of the same checked sign rules, so the preceding induction applies to each line of their trace.  Thus (15), strict positivity of \(\Delta\), and \(\varepsilon>0\) imply
\[
\llbracket H\rrbracket(z)\ge\varepsilon>0
\tag{17}
\]
for every point of the certificate's declared fiber.  In particular, substituting the response coordinates induced by the arbitrary \(\mathbf M\in\mathcal C(\mathcal P)\) proves the asserted context positivity.  A failed sign trace, \(\varepsilon\le0\), or a missing denominator certificate is rejected.  This proves validity of every accepted positivity certificate and makes the use of (14) noncircular: certificate soundness is an induction on the certificate trace, whereas (10) is an induction on the extraction derivation.

It remains to prove the claim about \(\operatorname{SupportLocal}_D\).  Put
\[
\Omega_{-D}=\prod_{D'\in\mathcal D(G),\,D'\ne D}\Omega_{D'}.
\]
All factors are finite and nonempty by \(\mathrm{def:district\text{-}response\text{-}profile}\); when there is no nonfocal district, the empty product is the singleton.  By \(\mathrm{ass:recursive\text{-}compatibility}\), acyclic substitution terminates and defines the Boolean table
\[
i_A(\omega_D,\omega_{-D})
=I_A(\omega_D,\omega_{-D})
\quad\text{on }\Omega_D\times\Omega_{-D}.
\tag{18}
\]
The checker enumerates this entire finite table.  It accepts exactly when
\[
\forall\omega_D\in\Omega_D\ \forall\omega_{-D},\omega'_{-D}\in\Omega_{-D},
\qquad
i_A(\omega_D,\omega_{-D})
=i_A(\omega_D,\omega'_{-D}).
\tag{19}
\]
Condition (19) is precisely pointwise independence of the recursive indicator from the nonfocal response profiles.  Thus the finite enumeration terminates and accepts if and only if the claimed independence holds.  On acceptance, let \(\omega_{-D}^0\) be the first profile in the fixed enumeration and recall that the constructed row is
\[
\operatorname{row}_D(A)(\omega_D)
=i_A(\omega_D,\omega_{-D}^0).
\tag{20}
\]
Equations (19)--(20) show, for every nonfocal profile, that
\[
\operatorname{row}_D(A)(\omega_D)
=i_A(\omega_D,\omega_{-D}).
\tag{21}
\]
Conversely, if (19) fails, exhaustive enumeration returns some \(\omega_D\) and two nonfocal profiles with unequal Boolean values; this is an explicit rejection witness and proves that no accepted row could equal the indicator uniformly.  Equations (10), (17), and (19)--(21) prove all assertions.

As a smallest exact check, take \(\Omega_D=\Omega_{-D}=\{0,1\}\).  The table \(i_A(\omega_D,\omega_{-D})=\omega_D\) passes and returns the row \((0,1)\), while \(i_A(\omega_D,\omega_{-D})=\omega_D\mathbin{\mathrm{xor}}\omega_{-D}\) fails with the two nonfocal profiles \(0\) and \(1\) at either fixed focal value.  For the first table, the complement calculation gives probability \(1-q_D^s(1)\); conditioning on the sure event has denominator one, whereas conditioning on the empty event has no positivity certificate and is rejected.  This checks both branches of the locality test and the zero-denominator boundary case on the smallest two-block carrier.